\documentclass[a4paper,12pt]{article}

\usepackage[T2A]{fontenc}
\usepackage[utf8]{inputenc}
\usepackage[russian]{babel}

\usepackage{amsmath, amssymb, amsthm, mathtools}
\usepackage{helvet}
\renewcommand{\familydefault}{\sfdefault}
\usepackage{icomma}
\usepackage{geometry}
\usepackage{microtype}
\usepackage{tikz}
\usepackage{tkz-euclide}
\usepackage{pgfplots}
\pgfplotsset{compat=1.18}
\usepackage{tabularx, booktabs, longtable}
\usepackage{enumitem}
\usepackage{hyperref}
\usepackage{parskip}
\usepackage{fancyhdr}
\usepackage{setspace}
\usepackage{xcolor}

\geometry{
  a4paper,
  left=18mm,
  right=18mm,
  top=18mm,
  bottom=20mm
}

\definecolor{accent}{HTML}{008080}
\definecolor{darkaccent}{HTML}{004D4D}
\definecolor{textgray}{HTML}{333333}
\definecolor{softgray}{HTML}{6E7777}

\hypersetup{
  colorlinks=true,
  linkcolor=darkaccent,
  urlcolor=darkaccent
}

\onehalfspacing
\setlength{\parindent}{0pt}
\setlist[itemize]{leftmargin=1.25em, itemsep=0.25em, topsep=0.35em}
\setlist[enumerate]{leftmargin=1.45em, itemsep=0.35em, topsep=0.35em}

\pagestyle{fancy}
\fancyhf{}
\lhead{\textcolor{textgray}{Метод интервалов}}
\rhead{\textcolor{textgray}{10 класс}}
\cfoot{\textcolor{textgray}{\thepage}}
\renewcommand{\headrulewidth}{0.4pt}
\renewcommand{\headrule}{\hbox to\headwidth{\color{accent}\leaders\hrule height \headrulewidth\hfill}}

\newcommand{\sectiontitle}[1]{%
  \vspace{0.65em}
  {\Large\bfseries\textcolor{darkaccent}{#1}}\par
  \vspace{0.25em}
  {\color{accent}\rule{\linewidth}{0.5pt}}
  \vspace{0.45em}
}

\newcommand{\subsectiontitle}[1]{%
  \vspace{0.5em}
  {\large\bfseries\textcolor{darkaccent}{#1}}\par
  \vspace{0.2em}
}

\newcommand{\important}[1]{\textbf{\textcolor{darkaccent}{#1}}}
\newcommand{\R}{\mathbb{R}}

\newcommand{\openpoint}[2]{%
  \draw[accent, line width=0.9pt, fill=white] (#1,0) circle (2.4pt)
    node[below=7pt, text=black] {$#2$};
}
\newcommand{\closedpoint}[2]{%
  \filldraw[accent] (#1,0) circle (2.4pt)
    node[below=7pt, text=black] {$#2$};
}

\begin{document}

\begin{titlepage}
\thispagestyle{empty}
\begin{center}
\vspace*{28mm}

{\Huge\bfseries\textcolor{darkaccent}{Чек-лист}}\par
\vspace{6mm}
{\LARGE\bfseries Метод интервалов: кратность корней и рациональные неравенства}\par

\vspace{18mm}

{\Large Лёвин Артём Александрович}\par
\vspace{2mm}
{\Large эксклюзивно для Тимофея Васильченко}\par

\vspace{12mm}
{\large 10 класс}\par
\vspace{4mm}
{\large \today}

\vfill

\begin{tikzpicture}[remember picture, overlay]
  \draw[accent, line width=1.1pt]
    ([xshift=18mm,yshift=24mm]current page.south west)
    .. controls ([xshift=55mm,yshift=48mm]current page.south west)
    and ([xshift=-65mm,yshift=8mm]current page.south east)
    .. ([xshift=-18mm,yshift=31mm]current page.south east);
  \draw[darkaccent, line width=0.6pt]
    ([xshift=30mm,yshift=17mm]current page.south west)
    .. controls ([xshift=72mm,yshift=35mm]current page.south west)
    and ([xshift=-92mm,yshift=46mm]current page.south east)
    .. ([xshift=-28mm,yshift=18mm]current page.south east);
\end{tikzpicture}

\end{center}
\end{titlepage}

\sectiontitle{1. Каркас метода интервалов}

Метод интервалов применяется к неравенствам, в которых слева стоит произведение или дробь, а справа --- ноль:
\[
F(x)\ge 0,\qquad F(x)>0,\qquad F(x)\le 0,\qquad F(x)<0.
\]

\important{Нормализованный вид:} все выражения перенесены в левую часть, справа стоит $0$, множители разложены настолько, насколько это возможно.

\subsectiontitle{Алгоритм}

\begin{enumerate}
  \item Приведите неравенство к нормализованному виду.
  \item Разложите числитель и знаменатель на множители.
  \item Найдите нули каждого множителя.
  \item Отдельно отметьте нули знаменателя: они запрещены при любом знаке неравенства.
  \item Расположите все критические точки на числовой прямой по возрастанию.
  \item Определите знак выражения на одном удобном интервале.
  \item При переходе через каждую точку учитывайте её кратность.
  \item Выберите интервалы с нужным знаком и расставьте скобки.
\end{enumerate}

\important{Калибровочное значение} выбирают внутри интервала, а затем подставляют в исходную левую часть. Значение выражения целиком вычислять необязательно: достаточно надёжно установить его знак.

\sectiontitle{2. Кратность корня и чередование знаков}

Если множитель имеет вид $(x-a)^m$, то $a$ --- корень кратности $m$.

\begin{center}
\renewcommand{\arraystretch}{1.4}
\begin{tabularx}{\linewidth}{@{}>{\bfseries}l X X@{}}
\toprule
Кратность & Что происходит со знаком & Как двигаться по оси \\
\midrule
Нечётная: $1,3,5,\ldots$ & Знак выражения меняется & При переходе через точку меняем $+$ на $-$ или $-$ на $+$ \\
Чётная: $2,4,6,\ldots$ & Знак выражения сохраняется & При переходе через точку оставляем прежний знак \\
\bottomrule
\end{tabularx}
\end{center}

\important{Короткое правило:} нечётная кратность --- знак меняется; чётная кратность --- знак сохраняется.

\subsectiontitle{Пример 1. Чётный и нечётный корни}

Решим неравенство
\[
(2x+1)^2(x-1)^3\ge 0.
\]

Критические точки:
\[
2x+1=0 \Rightarrow x=-\frac12 \quad (\text{кратность }2),
\qquad
x-1=0 \Rightarrow x=1 \quad (\text{кратность }3).
\]

На интервале $\left(-\frac12;1\right)$ удобно взять $x=0$:
\[
(2\cdot0+1)^2(0-1)^3=-1<0.
\]

Через $x=1$ знак меняется, через $x=-\frac12$ сохраняется:

\begin{center}
\begin{tikzpicture}[x=2.15cm,y=1cm]
  \draw[->,line width=0.7pt] (-2.25,0)--(2.35,0) node[right] {$x$};
  \closedpoint{-0.75}{-\frac12}
  \closedpoint{1.05}{1}
  \node[darkaccent,above=7pt] at (-1.5,0) {$-$};
  \node[darkaccent,above=7pt] at (0.15,0) {$-$};
  \node[darkaccent,above=7pt] at (1.7,0) {$+$};
  \node[softgray,above=24pt] at (-0.75,0) {чётная};
  \node[softgray,above=24pt] at (1.05,0) {нечётная};
\end{tikzpicture}
\end{center}

Поскольку неравенство нестрогое, оба допустимых нуля включаются:
\[
\boxed{x\in\left\{-\frac12\right\}\cup[1;+\infty)}.
\]

\subsectiontitle{Как меняется ответ при другом знаке}

Для того же выражения получаем:

\begin{center}
\renewcommand{\arraystretch}{1.4}
\begin{tabularx}{\linewidth}{@{}c X@{}}
\toprule
Неравенство & Ответ \\
\midrule
$(2x+1)^2(x-1)^3\ge0$ & $\left\{-\frac12\right\}\cup[1;+\infty)$ \\
$(2x+1)^2(x-1)^3>0$ & $(1;+\infty)$ \\
$(2x+1)^2(x-1)^3\le0$ & $(-\infty;1]$ \\
$(2x+1)^2(x-1)^3<0$ & $(-\infty;-\frac12)\cup\left(-\frac12;1\right)$ \\
\bottomrule
\end{tabularx}
\end{center}

\important{Будьте внимательны:} при строгом неравенстве корень чётной кратности исчезает из ответа и может разрезать один знакопостоянный промежуток на два.

\sectiontitle{3. Точки, которых знак не заставляет чередоваться}

Рассмотрим выражение
\[
(x+3)^{10}(4x-3)^{11}.
\]

Точка $x=-3$ имеет чётную кратность $10$, а точка $x=\frac34$ --- нечётную кратность $11$. При $x=0$ выражение отрицательно. Поэтому знак не меняется при переходе через $-3$ и меняется при переходе через $\frac34$.

Например,
\[
(x+3)^{10}(4x-3)^{11}\ge0
\quad\Longrightarrow\quad
x\in\{-3\}\cup\left[\frac34;+\infty\right).
\]

Если заменить показатель $11$ на $10$, обе кратности станут чётными. Тогда произведение неотрицательно при всех $x$:
\[
(x+3)^{10}(4x-3)^{10}\ge0
\quad\Longrightarrow\quad
x\in\R.
\]

\important{Вывод:} правило механического чередования знаков работает только для точек нечётной кратности. Кратность каждого корня нужно определять до расстановки знаков.

\sectiontitle{4. Рациональные неравенства}

Для дроби
\[
\frac{P(x)}{Q(x)}\mathrel{\gtreqless}0
\]
критические точки дают и числитель, и знаменатель:
\[
P(x)=0,\qquad Q(x)=0.
\]

\begin{center}
\renewcommand{\arraystretch}{1.45}
\begin{tabularx}{\linewidth}{@{}>{\bfseries}l X X@{}}
\toprule
Источник точки & Можно ли включить & Отметка на оси \\
\midrule
Нуль числителя & Да, если знак $\ge$ или $\le$ & Закрашенная точка при нестрогом неравенстве \\
Нуль знаменателя & Никогда: дробь не определена & Всегда выколотая точка \\
\bottomrule
\end{tabularx}
\end{center}

\subsectiontitle{Пример 2. Простейшая дробь}

Решим:
\[
\frac{x+2}{x-3}\ge0.
\]

Числитель обращается в ноль при $x=-2$. Знаменатель обращается в ноль при $x=3$, поэтому
\[
x\ne3.
\]

При $x=0$ дробь отрицательна. Обе критические точки имеют нечётную кратность, поэтому знак при переходе через каждую из них меняется:

\begin{center}
\begin{tikzpicture}[x=1.55cm,y=1cm]
  \draw[->,line width=0.7pt] (-3.0,0)--(3.2,0) node[right] {$x$};
  \closedpoint{-1.15}{-2}
  \openpoint{1.35}{3}
  \node[darkaccent,above=7pt] at (-2.05,0) {$+$};
  \node[darkaccent,above=7pt] at (0.1,0) {$-$};
  \node[darkaccent,above=7pt] at (2.3,0) {$+$};
  \node[softgray,above=24pt] at (-1.15,0) {числитель};
  \node[softgray,above=24pt] at (1.35,0) {знаменатель};
\end{tikzpicture}
\end{center}

Ответ:
\[
\boxed{x\in(-\infty;-2]\cup(3;+\infty)}.
\]

\sectiontitle{5. Особые случаи и контроль ошибок}

\subsectiontitle{5.1. У числителя нет действительных корней}

Отсутствие нулей числителя не означает отсутствие решений. Например, $x^2+1>0$ при всех $x$, поэтому знак дроби
\[
\frac{x^2+1}{x(3x+5)}
\]
полностью определяется знаком знаменателя. Для неравенства
\[
\frac{x^2+1}{x(3x+5)}<0
\]
получаем
\[
\boxed{x\in\left(-\frac53;0\right)}.
\]

\subsectiontitle{5.2. Сокращение множителей}

При сокращении одинакового множителя исходное ограничение сохраняется. Например,
\[
\frac{(x+2)^2(x-1)}{(x+2)(x-3)}
=\frac{(x+2)(x-1)}{x-3},
\qquad x\ne-2,\quad x\ne3.
\]

Точка $x=-2$ остаётся выколотой, хотя после сокращения знаменатель в ней уже не обращается в ноль.

\subsectiontitle{5.3. Финальный чек-лист}

\begin{itemize}
  \item Справа стоит $0$.
  \item Все множители разложены и все критические точки найдены.
  \item Точки записаны по возрастанию.
  \item Для каждой точки определены источник и кратность.
  \item Нули знаменателя выколоты при любом знаке неравенства.
  \item При чётной кратности знак сохранён, при нечётной --- изменён.
  \item При строгом неравенстве все нули исключены.
  \item Ответ записан объединением промежутков в порядке слева направо.
\end{itemize}

\newpage
\sectiontitle{Тренировочный блок}

Решите неравенства методом интервалов. В ответе укажите объединение промежутков или конечное множество.

\subsectiontitle{Средний уровень}
\begin{enumerate}[label=\textbf{\arabic*.}]
  \item $(x-2)^2(x+1)^3\ge0$.
  \item $(3x+2)^4(x-5)^3<0$.
  \item $\dfrac{x+4}{x-1}\le0$.
\end{enumerate}

\subsectiontitle{Выше среднего}
\begin{enumerate}[label=\textbf{\arabic*.},start=4]
  \item $(x-3)^2(x+2)^5\le0$.
  \item $(2x-1)^6(x+3)^2>0$.
  \item $\dfrac{x^2-5x+6}{(x-1)(x+4)}\ge0$.
  \item $\dfrac{(x+1)^2(x-4)}{(x-2)^3}\le0$.
  \item $\dfrac{x^2+1}{x(3x+5)}<0$.
  \item $\dfrac{(x-1)^4(x+2)}{(x-3)^2}\ge0$.
\end{enumerate}

\subsectiontitle{Сложный уровень}
\begin{enumerate}[label=\textbf{\arabic*.},start=10]
  \item $\dfrac{(x-1)^3(x+2)^2}{(x+2)(x-3)^2(x+4)}\le0$.
\end{enumerate}

\newpage
\sectiontitle{Ответы к тренировочному блоку}

\begin{center}
\renewcommand{\arraystretch}{1.5}
\begin{tabularx}{\linewidth}{@{}>{\bfseries}c X@{}}
\toprule
№ & Ответ \\
\midrule
1 & $[-1;+\infty)$ \\
2 & $(-\infty;-\frac23)\cup(-\frac23;5)$ \\
3 & $[-4;1)$ \\
4 & $(-\infty;-2]\cup\{3\}$ \\
5 & $\R\setminus\left\{-3;\frac12\right\}$ \\
6 & $(-\infty;-4)\cup(1;2]\cup[3;+\infty)$ \\
7 & $\{-1\}\cup(2;4]$ \\
8 & $\left(-\frac53;0\right)$ \\
9 & $[-2;3)\cup(3;+\infty)$ \\
10 & $(-\infty;-4)\cup(-2;1]$ \\
\bottomrule
\end{tabularx}
\end{center}



\end{document}